\documentclass[11pt]{article}
\usepackage[utf8]{inputenc}
\usepackage[margin=1in]{geometry}
\usepackage{amsmath,amssymb}
\usepackage{booktabs}
\usepackage{graphicx}
\usepackage{xcolor}
\usepackage[colorlinks=true,linkcolor=blue,citecolor=blue,urlcolor=blue]{hyperref}
\usepackage{authblk}

% ============================================================================
% WORKSHOP PAPER 1 -- evaluation-methodology angle
% All numerical results are real measurements from the OpenDriveVLA-0.5B study
% (nuScenes-mini open-loop + NeuroNCAP closed-loop, scenes 0103+0796, real cameras). See companion
% RESULTS.md / FINDINGS.md in the code release.
% ============================================================================

\title{\textbf{Open-Loop L2 Is Blind: Planning Metrics Hide Quantization\\
Damage in Autonomous-Driving Vision--Language--Action Models}}
\author[1]{Justin Williams}
\author[1]{Kishor Datta Gupta}
\author[1]{Roy George}
\author[1]{Mrinmoy Sarkar}
\affil[1]{Clark Atlanta University \quad \texttt{justinwilliamstech693@gmail.com}}
\date{}

\begin{document}
\maketitle

\begin{abstract}
Post-training quantization is routinely reported ``lossless'' for autonomous-driving
vision--language--action (VLA) models using the standard open-loop nuScenes L2 metric. We
show this conclusion is an \emph{artifact of the metric}. Using OpenDriveVLA-0.5B (a Qwen2.5
planner on UniAD perception), we first reproduce the lossless result---FP16, INT8, and INT4
weight quantization of the language backbone all yield $\approx 0.33$\,m L2. We then run a
single controlled ablation: zeroing the ego-state and history text handed to the planner
inflates L2 from $0.33$ to $6.38$\,m ($19\times$). Open-loop planning on nuScenes is therefore
$\approx$ ego-motion extrapolation; the perception/planning pathway---and any quantization
damage to it---contributes little to the reported number, which is why model size and bit-width
barely move it. The practical implication is a warning: \emph{do not certify a quantized driving
policy on open-loop L2}. We outline why closed-loop testing is necessary to expose
quantization-induced failures that open-loop metrics structurally cannot see.
\end{abstract}

\section{Introduction}
Compressing driving VLA policies for on-board inference is attractive, and the literature
generally reports that 4-bit weight quantization is essentially free on open-loop nuScenes
planning. We argue the standard metric cannot support that claim. We make three points, each
backed by a measurement on OpenDriveVLA-0.5B: (i) weight quantization down to INT4 is indeed
lossless \emph{on open-loop L2}; (ii) open-loop L2 is dominated by ego-motion extrapolation, so
it is structurally insensitive to the perception/planning computation; (iii) therefore the
``lossless'' conclusion says nothing about deployment behavior.

\section{Setup}
OpenDriveVLA-0.5B couples a Qwen2.5-0.5B language/action head (hidden size 896, 24 decoder
layers) to UniAD stage-1 track+map perception used as a ``vision tower.'' We evaluate on the
162 nuScenes-mini $\cap$ validation samples under the ST-P3 and UniAD L2 protocols. Weight-only
post-training quantization (per-output-channel symmetric round-to-nearest, dequantized to the
compute dtype) is applied \emph{only} to the 168 Linear layers of the Qwen backbone (357.8\,M
parameters); perception, the scene/track/map projectors, embeddings and the output head remain
in floating point, so the measured effect is purely the language head's W$b$ loss.

\section{Results}
\subsection{INT4 is ``lossless'' on open-loop L2}
\begin{table}[h]\centering
\caption{Open-loop L2 (m, ST-P3 cumulative protocol; lower is better), 162 mini$\cap$val samples.
Weight-only quantization of the Qwen backbone.}
\begin{tabular}{lccc}
\toprule
Precision & L2 @1/2/3s avg & well-formed output & note \\
\midrule
FP16 (baseline) & 0.33 & 100\% & matches reported 0.5B result \\
INT8            & 0.33 & 100\% & lossless \\
INT4            & 0.33 & 100\% & lossless \\
INT3            & 1.91 & 92\%  & text generation begins to break \\
INT2            & 6.45 & 0\%   & dead (all malformed) \\
\bottomrule
\end{tabular}
\end{table}
Down to INT4 the planner is indistinguishable from FP16. The degradation at INT3/INT2 is a
\emph{generation-coherence} failure (malformed trajectory text), not graded planning loss.

\subsection{Why the metric cannot see quantization: ego-ablation}
\begin{table}[h]\centering
\caption{Ego-ablation: zeroing the ego-state and 2\,s history text in the prompt. Same model,
same perception.}
\begin{tabular}{lcc}
\toprule
Condition & L2 (m) & relative \\
\midrule
Full input (FP16)          & 0.33 & $1\times$ \\
Ego-state + history zeroed & 6.38 & $19\times$ \\
\bottomrule
\end{tabular}
\end{table}
Removing the handed ego kinematics collapses the plan by $19\times$, while the outputs remain
well-formed. The model is therefore extrapolating the ego trajectory it is given; perception and
the language head contribute little to open-loop L2. This single number explains the otherwise
puzzling observations that (a) 0.5B/3B/7B variants report the same L2 and (b) INT8/INT4 are free:
there is little perception-grounded reasoning in this metric to lose.

\section{Discussion}
Open-loop L2 measures how well a model copies a velocity vector forward, not how well it perceives
or plans. Quantization that subtly perturbs perception or planning---or that destabilizes
generation under distribution shift---is invisible to it. Certifying a quantized driving policy on
open-loop L2 thus provides false assurance. The necessary test is \emph{closed-loop}: only there
is the logged ego trajectory unavailable to coast on, so perception/planning, and quantization
damage to them, can manifest. A companion study shows exactly such a failure (per-channel INT4
collapsing generation closed-loop where group-wise INT4 does not), invisible to the L2 metric here.

\section{Limitations}
The ego-ablation and quantization sweep are on a single model (OpenDriveVLA-0.5B) and the
nuScenes-mini split. The claim is about the \emph{metric's} structural insensitivity, supported by
the $19\times$ ablation; it does not assert a specific quantization technique is unsafe (that is the
subject of the closed-loop companion work).

\section{Conclusion}
4-bit weight quantization of a driving-VLA language head is lossless on open-loop nuScenes L2---but
the metric is $\approx$ ego-extrapolation ($19\times$ ego-ablation), so ``lossless'' is not evidence
of safety. Quantized driving policies must be evaluated closed-loop.

\begin{thebibliography}{9}
\bibitem{opendrivevla} OpenDriveVLA: \emph{Towards End-to-end Autonomous Driving with Large
Vision--Language--Action Models}. 2025. (Replace with verified citation.)
\bibitem{uniad} Hu et al. \emph{Planning-oriented Autonomous Driving (UniAD)}. CVPR 2023.
\bibitem{egostatus} Zhai et al. \emph{Rethinking the Open-Loop Evaluation of End-to-End Autonomous
Driving (``Ego status is all you need'')}. 2023.
\bibitem{neuroncap} NeuroNCAP: \emph{Photorealistic Closed-loop Safety Testing for Autonomous
Driving}. ECCV 2024.
\end{thebibliography}
\end{document}
